\section{导读：数据为什么重新成为产业主线}

人工智能常被说成由数据、算力、算法三驾马车驱动。过去几年，算力和模型架构更容易被看见；但当大语言模型接近数据墙、机器人又面临数据荒漠时，数据重新成为决定产业格局的关键变量。谢晨这期访谈的价值在于，它不是讲单一数据集，而是把数据产业拆成来源、闭环、金字塔、定价、Recipe 和版图。

\begin{importantbox}{本期核心命题}
数据不是静态原料，而是持续闭环。谁能更低成本地获得真实失败、仿真场景、可验证反馈和可复用配方，谁就能更快训练出可落地的智能系统。
\end{importantbox}

\begin{knowledgebox}{视觉策略说明}
本视频是固定访谈画面，没有教学 slides、白板或产品演示。按本仓库播客标准，正文不重复插入人物帧；封面用于来源识别，正文用数据金字塔、闭环和产业图谱承载教学内容。
\end{knowledgebox}

\subsection{本章小结}

本期是数据产业综述：LLM 的问题是高质量数据边际收益下降，机器人的问题是物理世界数据稀缺且昂贵。两者都需要数据闭环和评测体系。

